% Compile with: pdflatex lecture_33.tex
% Requires: amsmath, amssymb, amsthm, mathtools, xcolor, mdframed,
%           enumitem, hyperref, pagecolor, microtype

\documentclass[11pt]{article}
\usepackage[margin=1in]{geometry}
\usepackage{amsmath, amssymb, amsthm}
\usepackage{mathtools}
\usepackage{xcolor}
\usepackage{mdframed}
\usepackage{enumitem}
\usepackage{hyperref}
\usepackage{pagecolor}
\usepackage{microtype}

% ---------- Colour palette ----------
\definecolor{cream}{HTML}{FFFDF5}
\definecolor{defblue}{HTML}{1A4A7A}
\definecolor{thmgreen}{HTML}{1A5C3A}
\definecolor{warnAmber}{HTML}{7A4A00}
\definecolor{dangerRed}{HTML}{7A1A1A}
\definecolor{boxbluebg}{HTML}{EBF3FB}
\definecolor{boxgreenbg}{HTML}{EBF7EF}
\definecolor{boxamberbg}{HTML}{FDF5E6}
\definecolor{boxredbg}{HTML}{FBEBEB}
\pagecolor{cream}

% ---------- Theorem styles ----------
\newtheoremstyle{defstyle}    {8pt}{8pt}{\normalfont}{0pt}{\bfseries\color{defblue}} {.}{0.5em}{}
\newtheoremstyle{thmstyle}    {8pt}{8pt}{\normalfont}{0pt}{\bfseries\color{thmgreen}}{.}{0.5em}{}
\newtheoremstyle{notstyle}    {8pt}{8pt}{\normalfont}{0pt}{\bfseries\color{warnAmber}}{.}{0.5em}{}
\newtheoremstyle{cautionstyle}{8pt}{8pt}{\normalfont}{0pt}{\bfseries\color{dangerRed}}{.}{0.5em}{}

\theoremstyle{defstyle}
\newtheorem{definition}{Definition}[section]

\theoremstyle{thmstyle}
\newtheorem{theorem}{Theorem}[section]
\newtheorem{corollary}[theorem]{Corollary}
\newtheorem{lemma}[theorem]{Lemma}

\theoremstyle{notstyle}
\newtheorem*{remark}{Remark}
\newtheorem*{example}{Example}

\theoremstyle{cautionstyle}
\newtheorem*{caution}{Caution}

% ---------- Boxes ----------
\newmdenv[
  backgroundcolor=boxbluebg, linecolor=defblue, linewidth=2pt,
  topline=false, rightline=false, bottomline=false,
  innerleftmargin=10pt, innerrightmargin=10pt,
  innertopmargin=6pt, innerbottommargin=6pt,
  skipabove=8pt, skipbelow=8pt
]{defbox}

\newmdenv[
  backgroundcolor=boxgreenbg, linecolor=thmgreen, linewidth=2pt,
  topline=false, rightline=false, bottomline=false,
  innerleftmargin=10pt, innerrightmargin=10pt,
  innertopmargin=6pt, innerbottommargin=6pt,
  skipabove=8pt, skipbelow=8pt
]{thmbox}

\newmdenv[
  backgroundcolor=boxamberbg, linecolor=warnAmber, linewidth=2pt,
  topline=false, rightline=false, bottomline=false,
  innerleftmargin=10pt, innerrightmargin=10pt,
  innertopmargin=6pt, innerbottommargin=6pt,
  skipabove=8pt, skipbelow=8pt
]{notebox}

\newmdenv[
  backgroundcolor=boxredbg, linecolor=dangerRed, linewidth=2pt,
  topline=false, rightline=false, bottomline=false,
  innerleftmargin=10pt, innerrightmargin=10pt,
  innertopmargin=6pt, innerbottommargin=6pt,
  skipabove=8pt, skipbelow=8pt
]{cautionbox}

% ---------- Macros ----------
\newcommand{\Z}{\mathbb{Z}}
\newcommand{\N}{\mathbb{N}}
\newcommand{\R}{\mathbb{R}}
\newcommand{\Nd}{\mathbb{N}_{\text{d}}}

\begin{document}

\begin{center}
  {\Large\bfseries Lecture 33 --- Infinity of $\N$, $|\N| = |\N \times \N|$, and Cantor's Theorem}\\[4pt]
  {\normalsize CS 173: Discrete Structures \quad$\cdot$\quad Cardinality \& the Diagonal Argument}\\[2pt]
  {\small\color{gray} University of Illinois at Urbana-Champaign}
\end{center}
\vspace{1em}\hrule\vspace{1.5em}

\section{Setup and Goal}

We continue developing the theory of cardinality. Two large claims sit on the
horizon:
\[
  |\N| \;=\; |\Nd| \;=\; |\Z|, \qquad\qquad |\N| \;=\; |\N \times \N|.
\]
Before we can prove either, we need a foundational fact: the set $\N$ itself is
infinite. To formalise this we recall that, throughout this course, a natural
number $n \in \N$ is treated as the finite set $n = \{0, 1, \dots, n-1\}$, so a
statement like $|n| = |\N|$ is a statement about cardinalities of sets.

Our first target is therefore the following.

\begin{notebox}
\begin{remark}
We want to show
\[
  (\forall n \in \N)\bigl(|n| \neq |\N|\bigr),
\]
i.e.\ no finite natural number has the same cardinality as $\N$. Equivalently,
$\N$ is infinite.
\end{remark}
\end{notebox}

\section{$\N$ is infinite}

\begin{thmbox}
\begin{theorem}[$\N$ is infinite]
$\N$ is infinite. Equivalently,
\[
  (\forall n \in \N)\bigl(|n| \neq |\N|\bigr).
\]
\end{theorem}
\end{thmbox}

\begin{proof}
Towards a contradiction, suppose there exists $n \in \N$ such that $|n| = |\N|$.
Then there exists a bijection
\[
  f : n \longrightarrow \N.
\]
(Recall $n = \{0, 1, \dots, n-1\}$.)

\medskip
\noindent\textbf{Construct a witness that $f$ misses.} Define
\[
  S \;:=\; \sum_{i=0}^{n-1} f(i),
\]
and observe that $S \in \N$ because each $f(i) \in \N$ and $\N$ is closed under
finite sums.

\medskip
\noindent\textbf{Bound every output of $f$ by $S$.} For any
$i \in \{z \in \N \mid z < n\}$, we have
\[
  f(i)
  \;\leq\; f(i) + \!\!\sum_{\substack{j=0 \\ j \neq i}}^{n-1}\!\! f(j)
  \;=\; \sum_{j=0}^{n-1} f(j)
  \;=\; S.
\]
The first inequality holds because every term $f(j)$ being summed lies in
$\N$, hence is non-negative. Therefore
\[
  (\forall i \in n)\bigl(f(i) \;\leq\; S \;<\; S+1\bigr).
\]

\medskip
\noindent\textbf{Derive the contradiction.} Since $f$ is surjective and
$S+1 \in \N$, there exists $x \in \{z \in \N \mid z < n\}$ such that
\[
  f(x) \;=\; S + 1.
\]
However, by the bound just established, $f(x) < S+1$, so $f(x) \neq S+1$.
Contradiction.

\medskip
Therefore no such bijection exists, and $\N$ is infinite.
\end{proof}

\begin{notebox}
\begin{remark}[Why this proof works]
The argument is a clean cardinality version of the diagonal style: assuming a
bijection from a finite set onto $\N$ lets us \emph{construct} a specific
natural number, namely $S+1$, that is strictly larger than every value the
bijection produces. Surjectivity then forces a contradiction.
\end{remark}
\end{notebox}

\section{$|\N| = |\N \times \N|$}

Having established that $\N$ is genuinely infinite, we now compare it against
the seemingly larger set $\N \times \N$ of ordered pairs of naturals. Despite
appearances, the two have the same cardinality.

\begin{thmbox}
\begin{theorem}
$|\N| \;=\; |\N \times \N|.$
\end{theorem}
\end{thmbox}

\begin{proof}[Proof (in progress)]
By the Schr\"oder--Bernstein style strategy, it suffices to exhibit two
injections, one in each direction.

\medskip
\noindent\textbf{An injection $f : \N \to \N \times \N$.} Define
\[
  f(n) \;:=\; (0, n).
\]
This is plainly injective: if $f(n_1) = f(n_2)$, then $(0, n_1) = (0, n_2)$,
so $n_1 = n_2$.

\medskip
\noindent\textbf{An injection $g : \N \times \N \to \N$.} Define
\[
  g(x, y) \;:=\; 2^{x}\,\cdot\,3^{y}.
\]
We claim $g$ is injective. Suppose $g(x_1, y_1) = g(x_2, y_2)$, i.e.
\[
  2^{x_1} \cdot 3^{y_1} \;=\; 2^{x_2} \cdot 3^{y_2}.
\]
By the Fundamental Theorem of Arithmetic, the prime factorisation of a positive
integer is unique. Since $2$ and $3$ are distinct primes, the exponents on
each side must agree, giving $x_1 = x_2$ and $y_1 = y_2$. Hence
$(x_1, y_1) = (x_2, y_2)$, so $g$ is injective.

\medskip
With injections in both directions established, the conclusion
$|\N| = |\N \times \N|$ follows from the cardinal comparability theorem to be
applied next lecture. \qedhere
\end{proof}

\begin{notebox}
\begin{remark}[Why $2^{x} \cdot 3^{y}$?]
The choice of two distinct primes is what makes the encoding injective. Any two
distinct primes $p, q$ would do --- for example $g(x, y) = p^{x} q^{y}$ ---
because unique factorisation guarantees that recovering $(x, y)$ from
$g(x, y)$ is unambiguous. Picking $2$ and $3$ is just the smallest, cleanest
choice.
\end{remark}
\end{notebox}

\begin{cautionbox}
\begin{caution}
Avoid the temptation to write $g(x, y) = 2^{x} + 3^{y}$ or
$g(x, y) = x + y$. Neither is injective: e.g.\ $1 + 2 = 2 + 1$ but
$(1, 2) \neq (2, 1)$. Multiplicativity together with distinct prime bases is
exactly what unique factorisation needs.
\end{caution}
\end{cautionbox}

\section{Cantor's Theorem}

So far we have stacked $\N, \Z, \N \times \N$ on top of each other and found
they all share the same cardinality. A natural question: is there \emph{any}
operation that genuinely produces a strictly larger set? Cantor's theorem
answers yes --- the power set always does.

\begin{thmbox}
\begin{theorem}[Cantor's Theorem]
$\forall X \bigl(|X| < |P(X)|\bigr).$
\end{theorem}
\end{thmbox}

\begin{proof}[Proof (the diagonal argument)]
Let $X$ be a set. Towards a contradiction, suppose $|X| \geq |P(X)|$. Then
there exists a function
\[
  \varphi : X \longrightarrow P(X)
\]
such that $\varphi$ is surjective. Define the \emph{anti-diagonal set}
\[
  \Delta \;:=\; \{\, x \in X \,\mid\, x \notin \varphi(x) \,\}.
\]
Observe that $\Delta \subseteq X$, so $\Delta \in P(X)$.

\medskip
Since $\varphi$ is surjective and $\Delta \in P(X)$, there exists
$\delta \in X$ such that
\[
  \varphi(\delta) \;=\; \Delta.
\]
We now ask whether $\delta \in \Delta$, and split into the two exhaustive
cases.

\medskip
\noindent\textbf{Case 1: $\delta \in \Delta$.}\;
By definition of $\Delta$, membership requires $\delta \notin \varphi(\delta)$.
Since $\varphi(\delta) = \Delta$, this gives $\delta \notin \Delta$,
contradicting the case assumption.

\medskip
\noindent\textbf{Case 2: $\delta \notin \Delta$.}\;
Then $\neg(\delta \notin \varphi(\delta))$ fails to hold for $\delta$, i.e.\
$\delta$ does \emph{not} satisfy the membership condition for $\Delta$. Since
the only way to fail $x \notin \varphi(x)$ is to have $x \in \varphi(x)$, we
get $\delta \in \varphi(\delta) = \Delta$, contradicting the case assumption.

\medskip
Both cases yield a contradiction, so no surjection $\varphi : X \to P(X)$
exists. Combined with the obvious injection $x \mapsto \{x\}$ from $X$ into
$P(X)$ (which gives $|X| \leq |P(X)|$), we conclude
\[
  |X| \;<\; |P(X)|. \qedhere
\]
\end{proof}

\begin{notebox}
\begin{remark}[The case split, restated cleanly]
The whole proof hinges on a single self-referential question:
\emph{is $\delta$ a member of the set $\varphi(\delta) = \Delta$?} The
definition of $\Delta$ makes the two possible answers logically swap:
\[
  \delta \in \Delta \iff \delta \notin \varphi(\delta) \iff \delta \notin \Delta.
\]
That biconditional is the contradiction. Cases 1 and 2 just walk through it
in each direction.
\end{remark}
\end{notebox}

\begin{notebox}
\begin{remark}[Connection to Russell's paradox]
The same self-referential trick produces Russell's paradox: take $X$ to be
``the set of all sets'' and let $\varphi$ be the identity, so $\Delta$
becomes the set of all sets that do not contain themselves. The
contradiction here is precisely why no such universal set can exist in ZFC.
Cantor's theorem can be read as a constructive cardinality version of the
same observation.
\end{remark}
\end{notebox}

\begin{notebox}
\begin{remark}[Why ``diagonal'' / ``anti-diagonalising''?]
Imagine listing the elements of $X$ down the rows of an infinite table, and
labelling the columns with the same elements. The cell at row $x$, column $y$
records whether $y \in \varphi(x)$. The set $\Delta$ is built by walking
\emph{down the diagonal} ($x = y$ in each row) and flipping the answer:
$x$ goes into $\Delta$ exactly when $x \notin \varphi(x)$. By construction,
$\Delta$ disagrees with $\varphi(x)$ at the entry $x$ for every row $x$, so it
cannot be any row of the table. The same trick powers Cantor's proof that
$|\R| > |\N|$, and Russell's paradox.
\end{remark}
\end{notebox}

\begin{notebox}
\begin{remark}[Reference]
The construction is sometimes called the \emph{anti-diagonalising set} in the
literature; nLab catalogues it under ``diagonal argument''.
\end{remark}
\end{notebox}

\section{Summary}

\begin{enumerate}[leftmargin=*]
  \item $\N$ is infinite: no finite natural number $n$ satisfies $|n| = |\N|$.
        The proof builds the witness $S+1$ from the sum
        $S = \sum_{i=0}^{n-1} f(i)$ and uses surjectivity to force a
        contradiction.
  \item Adding pairs does not enlarge cardinality: $|\N| = |\N \times \N|$.
        The forward injection $n \mapsto (0, n)$ is trivial; the reverse
        injection $(x, y) \mapsto 2^{x} \cdot 3^{y}$ relies on unique prime
        factorisation.
  \item Cardinality arguments at this level have a recipe: \emph{either}
        construct an explicit bijection, \emph{or} construct two injections
        and appeal to the cardinal comparability theorem.
  \item \textbf{Cantor's theorem:} $|X| < |P(X)|$ for every set $X$. The proof
        opens with the anti-diagonal set
        $\Delta = \{x \in X \mid x \notin \varphi(x)\}$, which by construction
        cannot be in the image of any surjection $\varphi : X \to P(X)$.
\end{enumerate}

\begin{notebox}
\begin{remark}[Looking ahead]
Next lecture will (i) close out the chain $|\N| = |\Nd| = |\Z|$ and (ii)
finish combining today's two injections into the formal $|\N| = |\N \times \N|$
via the cardinal comparability theorem. With Cantor's theorem now in hand, the
canonical application is $|\N| < |P(\N)| = |\R|$ --- the existence of
cardinalities strictly larger than $|\N|$, and the beginning of the
infinite hierarchy of infinities.
\end{remark}
\end{notebox}

\end{document}